\documentclass[12pt,a4paper]{article}

% ── Packages ──────────────────────────────────────────────────────────────────
\usepackage[utf8]{inputenc}
\usepackage[T1]{fontenc}
\usepackage{geometry}
\usepackage{graphicx}
\usepackage{hyperref}
\usepackage{setspace}
\usepackage{titlesec}
\usepackage{booktabs}
\usepackage{array}
\usepackage{parskip}
\usepackage{microtype}
\usepackage{lmodern}
\usepackage{fancyhdr}
\usepackage{caption}
\usepackage{enumitem}
\usepackage{xcolor}
\usepackage{tocloft}

% ── Page geometry ─────────────────────────────────────────────────────────────
\geometry{
  a4paper,
  left=30mm,
  right=25mm,
  top=25mm,
  bottom=25mm,
}

% ── Hyperref setup ────────────────────────────────────────────────────────────
\hypersetup{
  colorlinks=true,
  linkcolor=black,
  urlcolor=blue!60!black,
  citecolor=black,
  pdftitle={The Large Synoptic Survey Telescope (LSST)},
  pdfauthor={Kshitij Govil},
}

% ── Header / Footer ───────────────────────────────────────────────────────────
\pagestyle{fancy}
\fancyhf{}
\fancyhead[L]{\small\textit{The Large Synoptic Survey Telescope}}
\fancyhead[R]{\small\textit{Kshitij Govil}}
\fancyfoot[C]{\thepage}
\renewcommand{\headrulewidth}{0.4pt}

% ── Section formatting ────────────────────────────────────────────────────────
\titleformat{\section}
  {\large\bfseries}{\thesection}{1em}{}[\titlerule]
\titleformat{\subsection}
  {\normalsize\bfseries}{\thesubsection}{1em}{}
\titleformat{\subsubsection}
  {\normalsize\itshape}{\thesubsubsection}{1em}{}

% ── Line spacing ──────────────────────────────────────────────────────────────
\onehalfspacing

% ── Figure caption style ──────────────────────────────────────────────────────
\captionsetup{
  font=small,
  labelfont=bf,
  labelsep=period,
}

% ─────────────────────────────────────────────────────────────────────────────
\begin{document}

% ── Title page ────────────────────────────────────────────────────────────────
\begin{titlepage}
  \centering
  \vspace*{30mm}
  {\LARGE\bfseries The Large Synoptic Survey Telescope (LSST)\par}
  \vspace{8mm}
  \rule{\textwidth}{0.5pt}
  \vspace{8mm}
  {\large Kshitij Govil\par}
  \vspace{4mm}
  {\normalsize Varad Lifestyle Consulting\par}
  {\normalsize Vancouver~$\cdot$~Jaipur\par}
  \vspace{8mm}
  {\small\url{https://dummyneversleeps.com}~~$\cdot$~~
         \url{https://github.com/inertialee}\par}
  \vspace{16mm}
  \begin{quote}
    \itshape
    ``Space is for everybody. It's not just for a few people in science or math,
    or for a select group of astronauts. That's our new frontier out there,
    and it's everybody's business to know about space.''
    \begin{flushright}
      --- \textsc{Christa McAuliffe}\\
      {\small Astronaut, The Challenger Space Shuttle, December~6, 1985}
    \end{flushright}
  \end{quote}
  \vfill
  {\small York University~$\cdot$~2020}
\end{titlepage}

% ── Table of Contents ─────────────────────────────────────────────────────────
\tableofcontents
\newpage

% ─────────────────────────────────────────────────────────────────────────────
\section{Introduction}
% ─────────────────────────────────────────────────────────────────────────────

When Charles Messier and Pierre M\'{e}chain created the list of Messier objects
in the late 18th century, the first 110 objects to appear were listed out of
frustration---separately (or so we are told), as non-comet objects seemingly
unwanted in their hunt for comets. It boasts of some deep sky objects as well,
making it one of the most valuable catalogues of all time.

A little over 200 years later, we had moved well beyond naked-eye observations
and photographic plates to the first digital survey, i.e.\ the Sloan Digital
Sky Survey (the year 2000). About the same time (1996--2001), from the
conceptual hopes of a Dark Matter Telescope and the Fifth Decadal Report
(\textit{Astronomy and Astrophysics in the New Millennium}), the ``Large
Aperture Synoptic Survey Telescope''---or now, the Large Synoptic Survey
Telescope, the LSST---was proposed. The basic design and objectives were
comprehensively laid out even at this juncture:

\begin{quote}
  ``The Large-aperture Synoptic Survey Telescope (LSST) is a 6.5-m-class
  optical telescope designed to survey the visible sky every week down to a
  much fainter level than that reached by existing surveys. It will catalog
  90 percent of the near-Earth objects larger than 300~m and assess the threat
  they pose to life on Earth. It will find some 10,000 primitive objects in the
  Kuiper Belt, which contains a fossil record of the formation of the solar
  system. It will also contribute to the study of the structure of the universe
  by observing thousands of supernovae, both nearby and at large redshift, and
  by measuring the distribution of dark matter through gravitational lensing.
  All the data will be available through the National Virtual Observatory,
  providing access for astronomers and the public to very deep images of the
  changing night sky.''
\end{quote}

For a project that was meant to take the better part of the next two decades,
the statement above seems remarkably crisp in hindsight.

% ─────────────────────────────────────────────────────────────────────────────
\section{General Overview}
% ─────────────────────────────────────────────────────────────────────────────

\subsection{Home Base and Affiliations}

The LSST was formally initiated by the LSST Corporation (non-profit,
headquartered in Tucson, Arizona), where nearly 40 institutional members and
34 international contributors representing 23 countries came together for the
advance of science and astronomy. The University of Arizona and the University
of Washington are founding members; the member list includes Google Inc.\ as
one well-known commercial entity.\footnote{The full list of members can be
found at \url{https://www.lsstcorporation.org/node/3}.} Its partnership with
the Association of Universities for Research in Astronomy (AURA) under the
National Science Foundation (NSF, USA), and with the Stanford Linear
Accelerator Center (SLAC) under the Department of Energy (DOE, USA) for
operations and data activities, are primary mentions in the list of
associations.

Broadly, the point runners for various parts of the project are:

\begin{itemize}[leftmargin=*]
  \item AURA to oversee the site construction.
  \item The telescope provided by the National Optical Astronomy
        Observatory.\footnote{Ground-based observatory in the USA.}
  \item The data operations centre to be created by the National Centre for
        Super-Computing Applications.\footnote{Part of the University of
        Illinois, USA.}
  \item The camera designed and constructed by the SLAC National Accelerator.
\end{itemize}

The current LSST team is led by Physicist and Stanford University Professor
Steven Kahn, amongst more than 100 astronomers, physicists, and engineers
spanning the globe.

\subsection{Site Location}

The LSST site is being constructed at El Pe\~{n}\'{o}n Cerro Pach\'{o}n,
Coquimbo region in north-central Chile. The Cerro Pach\'{o}n ridge, in the
foothills of the Andes Mountains, will be home to the LSST summit facility---a
region also known for the Gemini-South and the SOAR\footnote{Southern
Astrophysical Research (SOAR) telescope, Cerro Pach\'{o}n, Chile.}
telescopes. The piece of land is owned by AURA and is located roughly 100
kilometres from the LSST base facility in the town of La Serena. The region is
known for its exceedingly dry climate, making it suitable for infrared
observations, with very stable air overhead promising excellent seeing.

The summit facility spreads across two mountain peaks, amidst a main and an
auxiliary telescope, which have been tested for their rock strength and erosion
resistance. Following a detailed fluid dynamic analysis and weather testing, the
main telescope naturally finds itself at the highest and largest peak, with the
auxiliary one on a smaller peak farther east downhill. The orientation of the
facility ``steps'' down into a saddle area south-east of the main telescope,
serving as the service and operations building.

\begin{figure}[htbp]
  \centering
  % Image: Summit facility and wind direction
  % Source: https://www.lsst.org/about/tel-site/summit
  % To add: \includegraphics[width=0.85\textwidth]{images/summit_facility.jpg}
  \fbox{\parbox{0.85\textwidth}{\centering\vspace{20pt}
    \textit{[Figure 1: Summit Facility and Wind Direction]}\\
    \small Source: \url{https://www.lsst.org/about/tel-site/summit}
  \vspace{20pt}}}
  \caption{Summit facility and wind direction.
           Courtesy: \url{https://www.lsst.org/about/tel-site/summit}}
  \label{fig:summit}
\end{figure}

\subsection{Site Description}

The structure of the main telescope includes the telescope pier, a lower
enclosure supporting a rotating dome, and a 3{,}000 square metre service and
operations building. The auxiliary telescope is housed in a nearby enclosure.

A well-connected building is meant to provide a protected environment to care
for and move the mirrors around, while keeping temperature-controlled areas
separated to house observatory equipment. Heat-generating parts are carefully
maintained far enough from the equipment and below the service level, serviced
by an 80-tonne lift to move the mirrors and camera between the telescope and
the maintenance level. Magnetic motors will be employed to facilitate fast,
smooth, and quiet slewing, with the challenge being that the telescope must
slew 3.5 degrees to an adjacent field and settle in 4--5 seconds (4 seconds
for the structure, and 1 second for the mirror---5 seconds total between each
exposure). The pier and telescope mount hence need to be very stiff, with an
altitude-over-azimuth mount made of steel, hydrostatic bearings on both axes,
and a pier (16~m diameter; 1.25~m thick walls). All this is mounted directly
to virgin bedrock.

\begin{figure}[htbp]
  \centering
  % Image: Telescope enclosure view
  % Source: https://www.lsst.org/about/tel-site/summit
  % To add: \includegraphics[width=0.85\textwidth]{images/telescope_enclosure.jpg}
  \fbox{\parbox{0.85\textwidth}{\centering\vspace{20pt}
    \textit{[Figure 2: Telescope Enclosure View]}\\
    \small Source: \url{https://www.lsst.org/about/tel-site/summit}
  \vspace{20pt}}}
  \caption{Telescope enclosure view.
           Courtesy: \url{https://www.lsst.org/about/tel-site/summit}}
  \label{fig:enclosure}
\end{figure}

A sophisticated LSST observatory software scheduler will ensure minimal human
intervention during observing, making the facility well automated under the
watch of one operator at the summit.

\subsection{Cost and Duration of Project}

The start of the LSST as a project can be dated back to 2001 (Fifth Decadal
Report, Astronomy and Astrophysics, 2001---Large Aperture Synoptic Survey
Telescope). It was conceived as early as the 1990s, with Tony Tyson of the
University of California (Davis) and Roger Angel (University of Arizona)
visualising a telescope to survey dark matter and time domain (Dark Matter
Telescope, 1996).

With \$473 million as the upper cap for the entire project,\footnote{Major
Research Equipment and Facility Construction (MREFC) funding at a not-to-exceed
cost of \$473M.} AURA being one of the first contributors, the LSSTC was formed
in 2003. The Department of Energy (DOE/SLAC) contributed \$168M of the upper
cap for the camera. Notable individual contributors include Charles Simonyi
(US\$20M) and Bill Gates (upwards of US\$10M). Mirror construction began in
2007 and polishing was completed in 2014. Construction of the telescope and
site began in the fourth quarter of FY2014. Engineering first light is planned
for the end of Q1 FY2021, with full science operations starting at the end of
FY2022. In all, this is a project that will have taken almost two decades since
its official inception, and roughly 27 years in all.

\begin{figure}[htbp]
  \centering
  % Image: LSST site under construction
  % Source: https://www.lsst.org/about/tel-site
  % To add: \includegraphics[width=0.85\textwidth]{images/lsst_construction.jpg}
  \fbox{\parbox{0.85\textwidth}{\centering\vspace{20pt}
    \textit{[Figure 3: The LSST Site Under Construction]}\\
    \small Source: \url{https://www.lsst.org/about/tel-site}
  \vspace{20pt}}}
  \caption{The LSST site under construction.
           Courtesy: \url{https://www.lsst.org/about/tel-site}}
  \label{fig:construction}
\end{figure}

% ─────────────────────────────────────────────────────────────────────────────
\section{Technical Details}
% ─────────────────────────────────────────────────────────────────────────────

\subsection{Telescope}

The LSST boasts a unique three-mirror, three-lens optical assembly, along with
the world's largest charge-coupled device (CCD) camera. Magnetic motors power
its movements, and it has adopted an active optics system enabling shape
modification of the different mirrors to ensure improved image quality. The
telescopic assembly is equipped with light baffles to tackle ambient light,
wind protection, and thermal controls (natural ventilation and day-time
cooling). Its design is unique in the 8-m class of telescopes, with a very
wide field of view of 3.5 degrees in diameter, or 9.6 square degrees. This,
coupled with a large aperture, makes for a large \'{e}tendue (``light-spread'')
more than roughly three times that of the best existing telescopes---Subaru
and PAN-STARRS.

\subsection{Optical Assembly}

The most preliminary design for the LSST optical assembly came from Roger
Angel, Director of the Steward Observatory Mirror Lab, University of Arizona,
and his collaborators, as the ``Dark Matter Telescope'' design. It is the
latest in a long line of optical assembly improvements directed at removing
two aberrations---spherical and coma---leaving only astigmatism to enable a
wider field of view. The LSST then, as an improvement, uses a
``three-mirror anastigmat'' to cancel astigmatism as well. This enables
sharper images over a very wide field of view, with somewhat compromised light
gathering power due to the secondary mirror blocking the telescope aperture.

\subsection{Optics}

The system consists of three non-spherical mirrors (hence the name
three-mirror anastigmat): the primary mirror at 8.4~m, the secondary at
3.5~m convex, and the 5.0~m tertiary mirror. It can view a 9.62 square-degree
patch of the sky, a 3.5-degree field of view covering a 64~cm diameter flat
focal plane.

\begin{figure}[htbp]
  \centering
  % Image: Three mirror assembly and camera-detector
  % Source: https://www.lsst.org/about/tel-site/optical_design
  % To add: \includegraphics[width=0.85\textwidth]{images/three_mirror_assembly.jpg}
  \fbox{\parbox{0.85\textwidth}{\centering\vspace{20pt}
    \textit{[Figure 5: The Three Mirror Assembly and Camera-Detector]}\\
    \small Source: \url{https://www.lsst.org/about/tel-site/optical_design}
  \vspace{20pt}}}
  \caption{The three mirror assembly and camera-detector.
           Courtesy: \url{https://www.lsst.org/about/tel-site/optical_design}}
  \label{fig:threemirror}
\end{figure}

\begin{figure}[htbp]
  \centering
  % Image: The M1M3 monolith
  % Source: https://www.lsst.org/about/tel-site/optical_design
  % To add: \includegraphics[width=0.85\textwidth]{images/m1m3_monolith.jpg}
  \fbox{\parbox{0.85\textwidth}{\centering\vspace{20pt}
    \textit{[Figure 6: The M1M3 Monolith]}\\
    \small Source: \url{https://www.lsst.org/about/tel-site/optical_design}
  \vspace{20pt}}}
  \caption{The M1M3 monolith.
           Courtesy: \url{https://www.lsst.org/about/tel-site/optical_design}}
  \label{fig:m1m3}
\end{figure}

The system is described as fast-optical (M1: f/1.8; M2: f/1.00; M3: f/0.83),
coupled with a large optical area and sensitive detectors to take very deep
images with a very large field of view. The length of the system is compact at
6.4~m from the vertex of the secondary to that of the tertiary. The primary
is highly annular, giving an outer clear aperture of 8.36~m and an inner
diameter of 5.12~m---effectively a collecting area equivalent to a 6.67~m
filled aperture. A hole in the secondary mirror holds the camera readout
electronics; a hole in the tertiary holds equipment for optical alignment.

The system is optimised to work between the ultraviolet and near-infrared
electromagnetic wavelength range, following the filter set $u, g, r, i, z, y$.
Image quality is better than 0.3 arcseconds for all bands (80\% encircled
energy of a point source), and better than 0.2 arcseconds for the $r, i, z, y$
bands. The atmosphere limits imaging to around 0.7 arcseconds.

The detector captures 63\% of an object's light at the centre out to a field
radius of 0.7 degrees, and 57\% at the full field radius of 1.75 degrees. The
\'{e}tendue (aperture area $\times$ field of view) is
319.5~m$^{2}$\,deg$^{2}$.

\subsection{Mirror Design}

The primary and tertiary mirrors (M1 and M3) form a single continuous surface
known as the M1M3 monolith, fabricated from a single substrate of Ohara E6
low-expansion glass. Since the two mirrors curve differently, a slight cusp
forms where they meet.

To keep the structure light and thermally responsive, a honeycomb structure---
stiff yet lightweight---supports the mirror surface rather than a solid backing.
The mirror was cast and finished over seven years (2008--2015) at the mirror
lab of the University of Arizona's Steward Observatory.

The 3.42~m secondary mirror (M2) is the largest convex mirror ever made. A
100~mm-thick solid blank of low-expansion glass was used. By removing 1{,}735~lbs
of glass to leave a 19~mm face sheet, the mirror weight was reduced to 680~kg,
and the entire assembly to 2{,}714~kg---52\% of its initial weight.

\subsection{Camera}

The LSST camera is the largest ever constructed. It is the size of a small
car: approximately 5.5~ft by 9.8~ft in dimensions, weighing about 2{,}800~kg.
It is a wide-field optical imager capable of viewing light between
0.3--1~$\mu$m (near-ultraviolet to near-infrared). A 3.5-degree field of view
is enabled using a 10-$\mu$m pixel setup capable of 0.2 arcsecond sampling,
optimising pixel sensitivity versus resolution.

A flat image surface of approximately 64~cm diameter is achieved by a mosaic
of 189 16-megapixel silicon detectors, arranged on 21 raft assemblies, for a
total of approximately 3.2 gigapixels. A filter changing mechanism and shutter
are placed in the middle of the telescope.

\begin{figure}[htbp]
  \centering
  % Image: Schematic of the LSST camera
  % Source: https://www.lsst.org/about/camera
  % To add: \includegraphics[width=0.85\textwidth]{images/lsst_camera.jpg}
  \fbox{\parbox{0.85\textwidth}{\centering\vspace{20pt}
    \textit{[Figure 7: Schematic of the LSST Camera]}\\
    \small Source: \url{https://www.lsst.org/about/camera}
  \vspace{20pt}}}
  \caption{A schematic of the LSST camera.
           Courtesy: \url{https://www.lsst.org/about/camera}}
  \label{fig:camera}
\end{figure}

\subsubsection{Focal Plane}

Charge-coupled devices (CCDs) are laid out in an array operating at
approximately $-100$~\textdegree{}C, housed in an evacuated cryostat. The
lens at the cryostat serves both as the light entrance for the detector and
the vacuum seal for the cryostat. Five filters, each 75~cm in diameter, can
be accommodated inside the ``filter carousel'' and swapped robotically in
2 minutes or less.

\begin{figure}[htbp]
  \centering
  % Image: Camera focal plane schematic
  % Source: https://www.lsst.org/about/camera/focalplane
  % To add: \includegraphics[width=0.85\textwidth]{images/focal_plane.jpg}
  \fbox{\parbox{0.85\textwidth}{\centering\vspace{20pt}
    \textit{[Figure 8: The Camera's Focal Plane]}\\
    \small Source: \url{https://www.lsst.org/about/camera/focalplane}
  \vspace{20pt}}}
  \caption{The camera's focal plane---schematic.
           Courtesy: \url{https://www.lsst.org/about/camera/focalplane}}
  \label{fig:focalplane}
\end{figure}

\subsubsection{Detector Design}

A total of 21 rafts hold 189 CCDs, with each raft comprising a $3\times3$
square of sensors. Each CCD is built up of $4096 \times 4097$ imaging
elements with 16 outputs, giving 144 channels per raft and 3{,}024 channels
across the entire focal plane array.

The rafts are mounted on towers that also hold front-end electronics for data
readout. High-density flexible cables carry data to back-end electronics where
they are converted to digital format. Cooling straps maintain each raft at
approximately $-100$~\textdegree{}C. Each raft functions as an autonomous
complete camera and is individually controlled by the Observatory Control
System. The data rate is 8~megapixels per second per CCD, for a total of
3.2~gigapixels per second.

\begin{figure}[htbp]
  \centering
  % Image: 21 rafts assembled
  % Source: https://www.lsst.org/about/camera/rafttower
  % To add: \includegraphics[width=0.85\textwidth]{images/raft_towers.jpg}
  \fbox{\parbox{0.85\textwidth}{\centering\vspace{20pt}
    \textit{[Figure 9: 21 Rafts Mounted on Towers]}\\
    \small Source: \url{https://www.lsst.org/about/camera/rafttower}
  \vspace{20pt}}}
  \caption{21 rafts, each mounted on towers with electronics and assembled
           into one unit. Courtesy: \url{https://www.lsst.org/about/camera/rafttower}}
  \label{fig:rafts}
\end{figure}

\subsubsection{Specifications at a Glance}

\textbf{Focal Plane}
\begin{itemize}[leftmargin=*]
  \item High quantum efficiency to 1{,}000~nm
  \item Thick silicon $>100$~nm
  \item Point spread function effective at $\ll 0.7$~arcseconds
  \item Fast f/1.2 focal ratio
  \item Wide field of view at 3{,}200~cm$^{2}$ focal plane
  \item High throughput: $>90$\% fill factor
  \item 2-second fast readout; 3{,}024 output ports; 150 connections per sensor
  \item Low read noise: $<10$~electrons
\end{itemize}

\textbf{Camera Data Rates}
\begin{itemize}[leftmargin=*]
  \item More than 3~GB/s peak raw data
  \item 3.2 billion pixels readout in 2 seconds with 15-second integration
  \item Dynamic range: 18 bits/pixel
  \item $>0.6$~GB/s average in the pipeline
  \item Approximately 15~TB of data per night
  \item 2~Tflops/s average; 9~Tflops/s peak
\end{itemize}

% ─────────────────────────────────────────────────────────────────────────────
\section{Scientific Details}
% ─────────────────────────────────────────────────────────────────────────────

\subsection{The Science Case}

The roots of this project go back to the 1990s, with the idea of the Dark
Matter Telescope (Tony Tyson and Roger Angel). The LSST's scientific goals are
profound:

\begin{itemize}[leftmargin=*]
  \item What is the mysterious dark energy driving the acceleration of cosmic
        expansion?
  \item What is dark matter, and how do its distribution and properties affect
        the formation of stars, galaxies, and larger structures?
  \item What is the formation history of the Milky Way?
  \item What is the nature of the outer regions of the solar system?
  \item Are there new exotic and explosive phenomena in the universe yet to be
        discovered?
\end{itemize}

The LSST adopts a simple conceptual approach: conduct a deep survey over a
very large area of the sky; repeat images of every part of the sky fast enough;
pick up multiple bands of the electromagnetic spectrum every few nights; and
continue for 10 years.

\subsection{The Nature of Dark Matter and Dark Energy}

Several billion galaxies will be probed using a variety of cross-checking
methods. Strong and weak gravitational lensing provides great insight into
dark matter; the LSST will find thousands of gravitational lenses of all sizes
and configurations. Specific investigations include:

\begin{itemize}[leftmargin=*]
  \item Galaxy correlations in 3-D space vs.\ cosmic epoch to evaluate
        expansion of the universe.
  \item Counts of dark matter clusters via weak gravitational lensing fused
        with optical data, as probes of dark energy.
  \item Supernovae for probing the cosmic era when dark energy becomes
        prominent.
  \item Simultaneous measurements of mass growth and space-time curvature, to
        distinguish between dark energy and modified gravity as the source of
        recent acceleration.
\end{itemize}

\subsection{Cataloguing the Solar System}

Several million moving objects will be monitored (up to 100 times more than
current availability), encompassing orbital, colour, and variability
information. Near-Earth asteroids through to main belt asteroids will be
observed, and up to 90\% of Potentially Hazardous Asteroids (PHAs) larger
than 140~m in diameter could be detected.

\subsection{The Transient Sky}

Repeated night sky imaging at considerable depths with excellent image quality
will reveal new information about variable stars and cosmic explosions, while
also discovering entirely new classes of transient events. Within one minute
of a change in the sky, the LSST will generate an alert to the scientific
community. The LSST is forecast to discover 3--4 million supernovae during
its 10-year survey.

\subsection{Milky Way Structure and Formation}

The LSST will survey at least half the sky with each area observed over 100
times. Colours and brightness of billions of new stars, motions of millions
of them, and other observations will yield a map over 1{,}000 times deeper
than past surveys when stacked in depth. This is expected to answer fundamental
questions about the structure and evolutionary history of the Milky Way.

% ─────────────────────────────────────────────────────────────────────────────
\section{Special Techniques}
% ─────────────────────────────────────────────────────────────────────────────

\subsection{Survey Strategy}

The LSST is expected to conduct a 10-year survey taking more than five million
exposures, amounting to 50~petabytes of raw image data. Data will be taken in
back-to-back pairs of 15-second exposures. One pair constitutes a ``single
visit''---a single observation of a ten square-degree field through one
particular filter.

Operations strategies generated by the Operations Simulator (OpSim) include:

\begin{description}[leftmargin=*]
  \item[Wide, Fast, Deep (WFD)] 85--95\% of observing time; 18{,}000 square
        degrees of sky; declination range $-65^{\circ}$ to $+65^{\circ}$;
        pairs spaced $\sim$40 minutes apart across all six filters, observing
        each field every $\sim$3 days.
  \item[North Ecliptic Spur (NES)] Extension to WFD to reach higher airmass.
        Excludes $u$ and $y$ filters.
  \item[Galactic Plane] Regions of high stellar density. Fewer total
        exposures per field than WFD; no pairing.
  \item[South Celestial Pole (SCP)] Higher airmass extension to WFD, south
        of declination $-65^{\circ}$. Covers $ugrizy$.
  \item[Deep Drilling Fields (DD)] Single pointings in extended sequences;
        various filter combinations providing near-simultaneous colour
        information for variable and transient objects.
\end{description}

\subsection{Information Processing}

\begin{description}[leftmargin=*]
  \item[AlertSim] Generates alerts for sky changes and transmits them
        world-wide within 60 seconds of taking the exposure.
  \item[CatSim] Simulates properties and distributions of stars, galaxies,
        and asteroids.
  \item[PhoSim] Photon simulator; generates representative images of the sky.
  \item[GalSim] Open-source package for simulating images of stars and
        galaxies for weak-lensing analysis.
  \item[MAF] Metrics Analysis Framework; analyses outputs of the operations
        simulator.
\end{description}

% ─────────────────────────────────────────────────────────────────────────────
\section{Data Management}
% ─────────────────────────────────────────────────────────────────────────────

The LSST will produce about 20~TB of raw data per night. Data management is
architected in three layers:

\begin{enumerate}[leftmargin=*]
  \item \textbf{Infrastructure Layer:} Computing, storage, networking hardware
        and system software.
  \item \textbf{Middleware Layer:} Distributed processing, data access, user
        interface, and system operations services.
  \item \textbf{Applications Layer:} Data pipelines, data products, and
        science data archives.
\end{enumerate}

\subsection{Pipelines}

\begin{description}[leftmargin=*]
  \item[Nightly Pipelines] Based on image subtraction; highlight differences
        between two exposures of the same field. Detect transient events and
        send alerts within 60 seconds.
  \item[Data Release Pipelines] Release fully analysed data products covering
        long time scales.
  \item[Calibration Products Pipeline] Produces calibration data required by
        other pipelines.
\end{description}

\begin{figure}[htbp]
  \centering
  % Image: Data management pipelines
  % Source: https://www.lsst.org/about/dm
  % To add: \includegraphics[width=0.85\textwidth]{images/data_pipelines.jpg}
  \fbox{\parbox{0.85\textwidth}{\centering\vspace{20pt}
    \textit{[Figure 10: Data Management Pipelines]}\\
    \small Source: \url{https://www.lsst.org/about/dm}
  \vspace{20pt}}}
  \caption{Data management pipelines.
           Courtesy: \url{https://www.lsst.org/about/dm}}
  \label{fig:pipelines}
\end{figure}

\subsection{Data Products}

Three main categories:

\begin{description}[leftmargin=*]
  \item[Prompt] Generated every night with 60-second latency; source
        catalogues derived from difference images; image data products
        released with 24-hour latency.
  \item[Data Releases] Made available annually as the result of coherent
        processing of all science data to date.
  \item[User Generated] Originating from the community; access controls
        enable limited or open distribution.
\end{description}

The scientific database will include a source catalogue (7 trillion rows),
object catalogue (37 billion rows, 200 attributes each), a moving object
catalogue (6 million rows), a 60-second alerts database, and calibration,
configuration, processing, and provenance metadata.

\subsection{Data Storage and Processing Facilities}

\begin{itemize}[leftmargin=*]
  \item A mountain summit and base facility: data acquisition at the summit;
        storage and routing at the base.
  \item A supercomputing-class archive facility: complete processing,
        reprocessing, and permanent storage. IN2P3 (France) hosts a full
        copy and handles half the data release processing.
  \item Two project-funded data access centres co-located with the archive
        and base facilities; serving disaster recovery and open public
        interfaces to data products.
\end{itemize}

% ─────────────────────────────────────────────────────────────────────────────
\section{Beyond Surveying: Outreach and Education}
% ─────────────────────────────────────────────────────────────────────────────

While predominantly a survey project, the LSST aims to revolutionise public
and school education and outreach, reaching out under four categories:

\begin{description}[leftmargin=*]
  \item[General Public] An Education and Public Outreach (EPO) website
        featuring news about LSST discoveries, profiles of scientists and
        engineers, and opportunities for non-specialists to explore LSST data.
  \item[Formal Educators] An online data-driven classroom covering commonly
        taught topics in physics and astronomy.
  \item[Planetariums and Informal Science Centres] A gallery of high-end
        multimedia visualisations available for download and integration into
        exhibits and presentations.
  \item[Citizen Scientist Researchers] Support for citizen science projects
        using LSST data, including a dedicated project-building tool.
\end{description}

% ─────────────────────────────────────────────────────────────────────────────
% References
% ─────────────────────────────────────────────────────────────────────────────
\newpage
\section*{References}
\addcontentsline{toc}{section}{References}

\begin{enumerate}[leftmargin=*,label={[\arabic*]}]
  \item LSST Corporation. \textit{LSST: Large Synoptic Survey Telescope}.
        \url{https://www.lsst.org}. Accessed 2020.
  \item National Research Council. \textit{Astronomy and Astrophysics in the
        New Millennium}. National Academies Press, 2001.
  \item Tyson, J.~A. \textit{Large Synoptic Survey Telescope: Overview}.
        SPIE~4836, 2002.
  \item LSST Science Collaboration. \textit{LSST Science Book}, version 2.0.
        arXiv:0912.0201, 2009.
  \item LSST Camera Team, SLAC National Accelerator Laboratory.
        \url{https://www.lsst.org/about/camera}. Accessed 2020.
  \item LSST Data Management Team.
        \url{https://www.lsst.org/about/dm}. Accessed 2020.
\end{enumerate}

\vfill
\begin{center}
  \small\textcolor{gray}{%
    Kshitij Govil~$\cdot$~Varad Lifestyle Consulting~$\cdot$~
    \url{https://dummyneversleeps.com}~$\cdot$~
    \url{https://github.com/inertialee}%
  }
\end{center}

\end{document}
